%% kapa: exam=1 round=30 year=2019 period=1 form=A subject=economics subject_ko=경제학원론 seq=12 number=52
%% source: Q-Net 감정평가사 기출문제 (cst00309) artlSeq=5205205 file=extracted/1008/5205205/2019년 제30회 감정평가사 1차시험 문제지/2019년 제30회 감정평가사 1차시험 문제지/2019년도 제30회 감정평가사 1차 1교시 A형.hwp
%% conversion: kordoc 4.24.51 -> tools/gen1.py (auto); review=proofread-sonnet; flags=none
\begin{question}{52}
기업 A의 생산함수는 ${ Q=L+3K }$이다. 생산량이 일정할 때, 기업 A의 한계기술대체율에 관한 설명으로 옳은 것은? (단, ${ Q }$는 생산량, ${ L }$은 노동량, ${ K }$는 자본량, ${ Q>0, L>0, K>0 }$이다.)
\begin{choices}
  \item 노동과 자본의 투입량과 관계없이 일정하다.
  \item 노동 투입량이 증가하면 한계기술대체율은 증가한다.
  \item 노동 투입량이 증가하면 한계기술대체율은 감소한다.
  \item 자본 투입량이 증가하면 한계기술대체율은 증가한다.
  \item 자본 투입량이 증가하면 한계기술대체율은 감소한다.
\end{choices}
\end{question}
\answer{1}
